\scanpage{106}

Reste à exprimer \inserted{$(\overline{\mathrm B}')$} (18), pour ceci considérons \deleted{\unsure{pour|la} condition \unsure{forme lin.|formes lin.} $a'_{-1}$ qui s'annule sur les $a_i$ \unsure{l'hom.|le module}} $L=\unsure{E|\mathcal E}/k(a_0,\ldots,a_n)$, et l'hom. induit par $E\unsure{\text{---}|\longrightarrow}L$, $F\longrightarrow L$, \deleted{cet hom} je dis que c'est un épi, i.e. que c'est un épi sur \inserted{chaque} fibre \deleted{\unsure{de $f$|point}}. On est ramené au cas d'un corps de base, et à prouver que dans ce cas on n'a pas $F\subset k(a_0,\ldots,a_n)$, \deleted{ou encore que \unsure{la fo|Or}}

Or $F\cap k(a_0,\ldots,a_{\unsure{n+1|n-1}})=(0)$, car le système de $n$ formes linéaires $a'_1,\ldots,a'_n$ sur $k(a_{\unsure{0|j}}\ldots a_n)$ a une matrice triangulaire, de déterminant le produit des $c_i=\langle a_i,a'_{i+1}\rangle$ $(0\leq i\leq n-1)$, donc $\in k^*$. Comme $k(a_0,\ldots,a_{\unsure{n|n-1}})$ est de codim 1 dans $k(a_0,\ldots,a_n)$, \deleted{\unsure{cela ex}} et $\dim F=2$, cela implique qu'on ne peut avoir $F\subset k(a_0,\ldots a_n)$.

Soit \deleted{$D'$} $\Delta=\operatorname{Ker}(F\unsure{\longrightarrow|\text{---}}L)$, qui est une deuxième droite de $F$. La condition \inserted{$(\overline{18})$} sur $s_0$ signifie que $s_0\notin\Delta$ en chaque fibre. Donc la conjonction de $(\overline{17})$ $(\overline{18})$ $(\overline{19})$ signifie que $s_0$ est une section de $F$ \inserted{(fibré vect. de rg 2)} qui n'est dans $D$ ni $\Delta$ en aucune fibre -- il en existe tjrs \unsure{loc.\,lmt|localement}, OK.

% archive sheet number: 105
